\section{Related Work}
\label{sec:related}

\subsection{Salient Object Detection}
Deep SOD is dominated by fully-supervised encoder--decoder
networks~\cite{borji2015salient,wang2021salientsurvey}.
PiCANet~\cite{picanet2018} learns pixel-wise contextual attention, while
BASNet~\cite{basnet2019} and F3Net~\cite{f3net2020} emphasise boundary quality
through hybrid losses and feedback refinement. LDF~\cite{ldf2020} decouples body
and edge learning, PSGLoss~\cite{psgloss2021} introduces a pyramidal structural
loss, and RCSB~\cite{rcsb2022} performs recursive contour--saliency blending for
sharp boundaries. Lightweight designs such as CSNet~\cite{csnet2020},
EDN~\cite{edn2022}, ADMNet~\cite{sod_admnet}, BDFGNet~\cite{sod_bdfgnet}, and
STNet~\cite{sod_stnet} target efficiency on resource-constrained hardware.
Transformer backbones brought further gains: VST~\cite{vst2021} casts SOD as
token-to-token prediction, and more recent works such as MENet~\cite{menet2023}
and DC-Net~\cite{sod_dcnet} push accuracy with multi-stage enhancement and
divide-and-conquer decoding. \rev{Closest to our setting, MDSAM~\cite{mdsam2024}
adapts SAM to SOD with multi-scale adapters and a detail-enhancement module,
without text prompts; we include it in the comparison of \cref{tab:sota}.} These
methods learn ``objectness'' purely from
saliency data and do not exploit the semantic priors available in vision--language
foundation models. \methodname{} is complementary: it transfers the open-vocabulary
knowledge of CLIP and the promptable segmentation ability of SAM into the binary
SOD setting. For scenes where depth is available, RGB-D SOD
methods~\cite{sod_s2ma,sod_bbsnet,sod_jldcf,sod_spnet,sod_spsn} fuse a dedicated
depth stream to disambiguate cluttered foregrounds; we compare against them on the
SIP benchmark using RGB input only.

\subsection{Foundation Models for Dense Prediction}
SAM~\cite{kirillov2023segment} is a promptable, class-agnostic segmenter trained
on billion-scale masks; given point, box, or mask prompts it produces high-quality
object masks, but it carries no semantics and cannot decide \emph{which} object is
salient~\cite{wu2023samsurvey}. CLIP~\cite{radford2021learning} provides the
missing semantic alignment through contrastive image--text pretraining, and its
dense features have become a standard source of open-vocabulary supervision.
A growing body of work couples the two. CAT-Seg~\cite{cho2024cat} performs
cost aggregation over a dense CLIP correlation volume for open-vocabulary
segmentation, and ESC-Net~\cite{lee2024escnet} routes such a correlation into
SAM's decoder via pseudo prompts, iteratively refining it with stacked
enhance-segment-correlate blocks. Other open-vocabulary segmenters extract dense
labels directly from CLIP~\cite{zhou2022maskclip}, adapt CLIP with mask
proposals~\cite{liang2023ovseg}, or attach lightweight side
adapters~\cite{xu2023sanlang}. All of these target multi-class recognition;
\methodname{} adapts the SAM--CLIP correlation paradigm to class-agnostic binary
SOD, \rev{and measures which of the adaptation choices this requires actually
matter.}
%% OLD: "is, to our knowledge, the first to adapt ... and it identifies and fixes the
%% gradient-flow issue that arises when the correlation is iteratively refined."

\subsection{Prompt Learning for Vision--Language Models}
Because CLIP's zero-shot behaviour is sensitive to prompt wording, a line of work
learns prompts rather than hand-crafting them: CoOp~\cite{zhou2022coop} optimises
continuous context vectors, and CoCoOp~\cite{zhou2022cocoop} conditions them on
the input image to improve generalisation to unseen classes. Our prompt bank
is related but operates differently: instead of learning the prompt tokens, we keep
a fixed bank of $K$ natural-language salient-object descriptions and \rev{study
how to \emph{aggregate} their correlation maps (max-pooling or learned per-pixel
attention). In our study (\cref{sec:study_prompt}) the bank size, the wording and
the aggregation mode have no consistent effect on S-measure once the visual
encoder is fine-tuned.}
%% OLD: "Our analysis shows that for SOD the aggregation mechanism matters more than
%% the specific wording (randomly sampled prompts with learned attention match
%% carefully curated prompts), which suggests that spatial aggregation, not
%% prompt-token optimisation, is the key lever ... This is a distinct finding from
%% the classification-oriented prompt-learning literature, where the exact context
%% tokens are decisive."
